% GENERATED FILE. Edit canonical lessons, then run npm run generate:books.
% canonical-source-hash: fnv1a64:9a1477e99c7fdb9b
% canonical-lessons: KA-C468-gatra, KA-C468-asali, KA-C468-gunamatta, KA-C468-khatari, KA-C468-masu, KA-R468-review-462-468

\chapter{Size and Quality}
\label{ch:ka-size-and-quality}
\hlchaptermodality{468}

\begin{chapteropening}
\textbf{By the end of this chapter:} \emph{I can say size, genuine, quality, a guarantee and to fade.}

Complete the last lesson of chapter 468: I can say size, genuine, quality, a guarantee and to fade.
\end{chapteropening}

\section[\texorpdfstring{\kn{ಗಾತ್ರ} (gātra)}{ಗಾತ್ರ (gātra)}]{\kn{ಗಾತ್ರ} (gātra) — size}

\label{lesson:KA-C468-gatra}



\begin{quote}
Before the new one: say the Kannada for a sale, then the Kannada for to buy for oneself.
\end{quote}

\subsection*{You'll want to know: \kn{ಗಾತ್ರ}}
\textbf{\kn{ಗಾತ್ರ}} — \emph{gātra} — \textquotedblleft{}size\textquotedblright{}.

\begin{quote}
\kn{ಈ} \kn{ಅಂಗಿ} \kn{ನನ್ನ} \kn{ಗಾತ್ರಕ್ಕೆ} \kn{ಸರಿಯಾಗಿಲ್ಲ}.
\end{quote}

\begin{quote}
\emph{ī aṅgi nanna gātrakke sariyāgilla.}
\end{quote}

\begin{quote}
This shirt is not my size.
\end{quote}

In Sanskrit a limb or the body; Kannada uses it for the size of things.

\begin{cousinweb}
\textbf{\kn{ಗಾತ್ರ}} (\emph{gātra}) is borrowed from Sanskrit \emph{gātra}, \textquotedblleft{}a limb, the body\textquotedblright{} (CDIAL 4124); in Kannada it came to mean bulk and size.
\end{cousinweb}

\begin{grammarlens}[title={What you've built}]
One more word for getting around.
\end{grammarlens}

\subsection*{Your turn}
\begin{itemize}
\raggedright
  \item \emph{Say it:} \emph{gātra}
  \item \emph{Say it:} \emph{gātra}, once more
  \item \emph{Say it:} \emph{ī aṅgi nanna gātrakke sariyāgilla.} — This shirt is not my size.
  \item \emph{Recall:} say the Kannada for a sale, then the Kannada for to buy for oneself, then say \emph{gātra} again
\end{itemize}

\begin{tcolorbox}[breakable,colback=teal!4,colframe=teal!35!black,title={Before you move on}]
What does \kn{ಗಾತ್ರ} mean? (\textquotedblleft{}Size\textquotedblright{}.) Say it once more.
\end{tcolorbox}
\section[\texorpdfstring{\kn{ಅಸಲಿ} (asali)}{ಅಸಲಿ (asali)}]{\kn{ಅಸಲಿ} (asali) — genuine; real}

\label{lesson:KA-C468-asali}



\begin{quote}
Before the new one: say the Kannada for to buy for oneself, then the Kannada for size.
\end{quote}

\subsection*{You'll want to know: \kn{ಅಸಲಿ}}
\textbf{\kn{ಅಸಲಿ}} — \emph{asali} — \textquotedblleft{}genuine; real\textquotedblright{}.

\begin{quote}
\kn{ಇದು} \kn{ಅಸಲಿ} \kn{ರೇಷ್ಮೆ} \kn{ಅಂತ} \kn{ಅಂಗಡಿಯವನು} \kn{ಹೇಳಿದನು}.
\end{quote}

\begin{quote}
\emph{idu asali rēṣme anta aṅgaḍiyavanu hēḷidanu.}
\end{quote}

\begin{quote}
The shopkeeper said this is genuine silk.
\end{quote}

Its opposite is \kn{ನಕಲಿ}, a fake, from \kn{ನಕಲು}, a copy.

\begin{cousinweb}
\textbf{\kn{ಅಸಲಿ}} (\emph{asali}) is usually traced to Arabic \emph{aṣl}, \textquotedblleft{}root, origin\textquotedblright{}, through Hindustani \emph{aslī}, \textquotedblleft{}original, genuine\textquotedblright{}.
\end{cousinweb}

\begin{grammarlens}[title={What you've built}]
One more word for getting around.
\end{grammarlens}

\subsection*{Your turn}
\begin{itemize}
\raggedright
  \item \emph{Say it:} \emph{asali}
  \item \emph{Say it:} \emph{asali}, once more
  \item \emph{Say it:} \emph{idu asali rēṣme anta aṅgaḍiyavanu hēḷidanu.} — The shopkeeper said this is genuine silk.
  \item \emph{Recall:} say the Kannada for to buy for oneself, then the Kannada for size, then say \emph{asali} again
\end{itemize}

\begin{tcolorbox}[breakable,colback=teal!4,colframe=teal!35!black,title={Before you move on}]
What does \kn{ಅಸಲಿ} mean? (\textquotedblleft{}Genuine; real\textquotedblright{}.) Say it once more.
\end{tcolorbox}
\section[\texorpdfstring{\kn{ಗುಣಮಟ್ಟ} (guṇamaṭṭa)}{ಗುಣಮಟ್ಟ (guṇamaṭṭa)}]{\kn{ಗುಣಮಟ್ಟ} (guṇamaṭṭa) — quality; standard}

\label{lesson:KA-C468-gunamatta}



\begin{quote}
Before the new one: say the Kannada for size, then the Kannada for genuine.
\end{quote}

\subsection*{You'll want to know: \kn{ಗುಣಮಟ್ಟ}}
\textbf{\kn{ಗುಣಮಟ್ಟ}} — \emph{guṇamaṭṭa} — \textquotedblleft{}quality; standard\textquotedblright{}.

\begin{quote}
\kn{ಈ} \kn{ಅಂಗಡಿಯ} \kn{ಬಟ್ಟೆಗಳ} \kn{ಗುಣಮಟ್ಟ} \kn{ಚೆನ್ನಾಗಿದೆ}.
\end{quote}

\begin{quote}
\emph{ī aṅgaḍiya baṭṭegaḷa guṇamaṭṭa cennāgide.}
\end{quote}

\begin{quote}
The quality of the clothes in this shop is good.
\end{quote}

\kn{ಗುಣ}, a quality, plus \kn{ಮಟ್ಟ}, a level.

\begin{cousinweb}
\textbf{\kn{ಗುಣಮಟ್ಟ}} (\emph{guṇamaṭṭa}) joins Sanskrit \emph{guṇa}, \textquotedblleft{}a quality\textquotedblright{} (CDIAL 4190), and native \emph{maṭṭa}, \textquotedblleft{}measure, level, limit\textquotedblright{} (DEDR 4660).
\end{cousinweb}

\begin{grammarlens}[title={What you've built}]
One more word for getting around.
\end{grammarlens}

\subsection*{Your turn}
\begin{itemize}
\raggedright
  \item \emph{Say it:} \emph{guṇamaṭṭa}
  \item \emph{Say it:} \emph{guṇamaṭṭa}, once more
  \item \emph{Say it:} \emph{ī aṅgaḍiya baṭṭegaḷa guṇamaṭṭa cennāgide.} — The quality of the clothes in this shop is good.
  \item \emph{Recall:} say the Kannada for size, then the Kannada for genuine, then say \emph{guṇamaṭṭa} again
\end{itemize}

\begin{tcolorbox}[breakable,colback=teal!4,colframe=teal!35!black,title={Before you move on}]
What does \kn{ಗುಣಮಟ್ಟ} mean? (\textquotedblleft{}Quality; standard\textquotedblright{}.) Say it once more.
\end{tcolorbox}
\section[\texorpdfstring{\kn{ಖಾತರಿ} (khātari)}{ಖಾತರಿ (khātari)}]{\kn{ಖಾತರಿ} (khātari) — a guarantee; an assurance}

\label{lesson:KA-C468-khatari}



\begin{quote}
Before the new one: say the Kannada for genuine, then the Kannada for quality.
\end{quote}

\subsection*{You'll want to know: \kn{ಖಾತರಿ}}
\textbf{\kn{ಖಾತರಿ}} — \emph{khātari} — \textquotedblleft{}a guarantee; an assurance\textquotedblright{}.

\begin{quote}
\kn{ಈ} \kn{ಫ್ಯಾನಿಗೆ} \kn{ಎರಡು} \kn{ವರ್ಷದ} \kn{ಖಾತರಿ} \kn{ಇದೆ}.
\end{quote}

\begin{quote}
\emph{ī phyānige eraḍu varṣada khātari ide.}
\end{quote}

\begin{quote}
This fan has a two-year guarantee.
\end{quote}

\kn{ಖಾತರಿ} \kn{ಕೊಡು}, to give an assurance.

\begin{cousinweb}
\textbf{\kn{ಖಾತರಿ}} (\emph{khātari}) came into Kannada through Hindustani; it is usually traced to Arabic \emph{khāṭir}, \textquotedblleft{}the mind, regard\textquotedblright{}, but how it came to mean an assurance is not settled.
\end{cousinweb}

\begin{grammarlens}[title={What you've built}]
One more word for getting around.
\end{grammarlens}

\subsection*{Your turn}
\begin{itemize}
\raggedright
  \item \emph{Say it:} \emph{khātari}
  \item \emph{Say it:} \emph{khātari}, once more
  \item \emph{Say it:} \emph{ī phyānige eraḍu varṣada khātari ide.} — This fan has a two-year guarantee.
  \item \emph{Recall:} say the Kannada for genuine, then the Kannada for quality, then say \emph{khātari} again
\end{itemize}

\begin{tcolorbox}[breakable,colback=teal!4,colframe=teal!35!black,title={Before you move on}]
What does \kn{ಖಾತರಿ} mean? (\textquotedblleft{}A guarantee; an assurance\textquotedblright{}.) Say it once more.
\end{tcolorbox}
\section[\texorpdfstring{\kn{ಮಾಸು} (māsu)}{ಮಾಸು (māsu)}]{\kn{ಮಾಸು} (māsu) — to fade; to grow dull}

\label{lesson:KA-C468-masu}



\begin{quote}
Before the new one: say the Kannada for quality, then the Kannada for a guarantee.
\end{quote}

\subsection*{You'll want to know: \kn{ಮಾಸು}}
\textbf{\kn{ಮಾಸು}} — \emph{māsu} — \textquotedblleft{}to fade; to grow dull\textquotedblright{}.

\begin{quote}
\kn{ಎರಡು} \kn{ಸಲ} \kn{ತೊಳೆದ} \kn{ಮೇಲೆ} \kn{ಅಂಗಿಯ} \kn{ಬಣ್ಣ} \kn{ಮಾಸಿತು}.
\end{quote}

\begin{quote}
\emph{eraḍu sala toḷeda mēle aṅgiya baṇṇa māsitu.}
\end{quote}

\begin{quote}
After two washes the colour of the shirt faded.
\end{quote}

Past (it …): \kn{ಮಾಸಿತು}, \emph{māsitu}.

Of colours, cloth and old paint.

\begin{cousinweb}
\textbf{\kn{ಮಾಸು}} (\emph{māsu}) is native Dravidian: DEDR 4781 lists Kannada \emph{māsu}, \textquotedblleft{}to grow dim, to be obscured, to fade\textquotedblright{}.
\end{cousinweb}

\begin{grammarlens}[title={What you've built}]
One more word for getting around.
\end{grammarlens}

\subsection*{Your turn}
\begin{itemize}
\raggedright
  \item \emph{Say it:} \emph{māsu}
  \item \emph{Say it:} \emph{māsu}, once more
  \item \emph{Say it:} \emph{eraḍu sala toḷeda mēle aṅgiya baṇṇa māsitu.} — After two washes the colour of the shirt faded.
  \item \emph{Recall:} say the Kannada for quality, then the Kannada for a guarantee, then say \emph{māsu} again
\end{itemize}

\begin{tcolorbox}[breakable,colback=teal!4,colframe=teal!35!black,title={Before you move on}]
What does \kn{ಮಾಸು} mean? (\textquotedblleft{}To fade; to grow dull\textquotedblright{}.) Say it once more.
\end{tcolorbox}
\section[(dialogue)]{Looking back: from In the Hotel Room to Size and Quality}

\label{lesson:KA-R468-review-462-468}



\begin{quote}
Say the Kannada for a balcony and an elevator.
\end{quote}

\subsection*{Your turn}
Say these aloud:

\begin{itemize}
\raggedright
  \item a balcony, an elevator, a mosquito net, a heater, to creak
  \item a service, a facility, to provide, a water heater, a resort
  \item breakfast, a menu, a waiter, a gratuity, to order
  \item a small town, a capital city, a region, a place, a centre
  \item a suburb, to reside, a town council, a streetlight, a drain
  \item a marketplace, a bazaar, a store, a sale, to buy for oneself
  \item size, genuine, quality, a guarantee, to fade
\end{itemize}

\begin{tcolorbox}[breakable,colback=teal!4,colframe=teal!35!black,title={Before you move on}]
A balcony? (\textbf{\kn{ಬಾಲ್ಕನಿ}}, \emph{bālkani}.) An elevator? (\textbf{\kn{ಲಿಫ್ಟ್}}, \emph{liphṭ}.) A mosquito net? (\textbf{\kn{ಸೊಳ್ಳೆಪರದೆ}}, \emph{soḷḷeparade}.) A heater? (\textbf{\kn{ಹೀಟರ್}}, \emph{hīṭar}.) To creak? (\textbf{\kn{ಕಿರುಗುಟ್ಟು}}, \emph{kiruguṭṭu}.) A service? (\textbf{\kn{ಸೇವೆ}}, \emph{sēve}.) A facility? (\textbf{\kn{ಸೌಲಭ್ಯ}}, \emph{saulabhya}.) To provide? (\textbf{\kn{ಒದಗಿಸು}}, \emph{odagisu}.) A water heater? (\textbf{\kn{ಗೀಸರ್}}, \emph{gīsar}.) A resort? (\textbf{\kn{ರೆಸಾರ್ಟ್}}, \emph{resārṭ}.) Breakfast? (\textbf{\kn{ಉಪಹಾರ}}, \emph{upahāra}.) A menu? (\textbf{\kn{ಮೆನು}}, \emph{menu}.) A waiter? (\textbf{\kn{ಪರಿಚಾರಕ}}, \emph{paricāraka}.) A gratuity? (\textbf{\kn{ಭಕ್ಷೀಸು}}, \emph{bhakṣīsu}.) To order? (\textbf{\kn{ಆರ್ಡರ್} \kn{ಮಾಡು}}, \emph{ārḍar māḍu}.) A small town? (\textbf{\kn{ಪಟ್ಟಣ}}, \emph{paṭṭaṇa}.) A capital city? (\textbf{\kn{ರಾಜಧಾನಿ}}, \emph{rājadhāni}.) A region? (\textbf{\kn{ಪ್ರದೇಶ}}, \emph{pradēśa}.) A place? (\textbf{\kn{ಸ್ಥಳ}}, \emph{sthaḷa}.) A centre? (\textbf{\kn{ಕೇಂದ್ರ}}, \emph{kēndra}.) A suburb? (\textbf{\kn{ಉಪನಗರ}}, \emph{upanagara}.) To reside? (\textbf{\kn{ವಾಸಿಸು}}, \emph{vāsisu}.) A town council? (\textbf{\kn{ಪುರಸಭೆ}}, \emph{purasabhe}.) A streetlight? (\textbf{\kn{ಬೀದಿದೀಪ}}, \emph{bīdidīpa}.) A drain? (\textbf{\kn{ಚರಂಡಿ}}, \emph{caraṇḍi}.) A marketplace? (\textbf{\kn{ಮಾರುಕಟ್ಟೆ}}, \emph{mārukaṭṭe}.) A bazaar? (\textbf{\kn{ಪೇಟೆ}}, \emph{pēṭe}.) A store? (\textbf{\kn{ಮಳಿಗೆ}}, \emph{maḷige}.) A sale? (\textbf{\kn{ಮಾರಾಟ}}, \emph{mārāṭa}.) To buy for oneself? (\textbf{\kn{ಕೊಂಡುಕೊ}}, \emph{koṇḍuko}.) Size? (\textbf{\kn{ಗಾತ್ರ}}, \emph{gātra}.) Genuine? (\textbf{\kn{ಅಸಲಿ}}, \emph{asali}.) Quality? (\textbf{\kn{ಗುಣಮಟ್ಟ}}, \emph{guṇamaṭṭa}.) A guarantee? (\textbf{\kn{ಖಾತರಿ}}, \emph{khātari}.) To fade? (\textbf{\kn{ಮಾಸು}}, \emph{māsu}.)
\end{tcolorbox}
